\documentclass[11pt]{article}
\usepackage{amsmath,amssymb,amsthm}

\newtheorem{theorem}{Theorem}

\begin{document}

\begin{theorem}[Pick's Theorem]
\label{thm:pick}
Let $P$ be a simple polygon in the plane whose vertices all have integer
coordinates. Let $i$ be the number of lattice points in the interior of $P$,
and let $b$ be the number of lattice points on the boundary of $P$ (including
the vertices). Then the area $A$ of $P$ satisfies
\[
A=i+\frac{b}{2}-1.
\]
\end{theorem}

\begin{proof}
Add noncrossing diagonals until $P$ is triangulated, using every lattice
point of $P$ as a vertex of the triangulation. Then each triangle in the
triangulation is a primitive lattice triangle, meaning that it has no lattice
points other than its three vertices on its boundary or in its interior. By
the standard lattice-triangle lemma, every primitive lattice triangle has area
$\tfrac12$.

Let $T$ be the number of triangles in the triangulation. Since area is
additive,
\[
A=\frac{T}{2}.
\]

Now view the triangulation as a planar graph. Let $V$, $E$, and $F$ be the
numbers of vertices, edges, and faces. Because every lattice point of $P$ is a
vertex, we have
\[
V=i+b.
\]
The faces are the $T$ bounded triangles together with the unbounded exterior
face, so
\[
F=T+1.
\]
Euler's formula gives
\[
V-E+F=2.
\]

To count edges, let $E_{\mathrm{int}}$ be the number of interior edges. Then
\[
E=E_{\mathrm{int}}+b,
\]
because the polygon boundary contributes exactly $b$ edges, one between each
pair of consecutive boundary lattice points. Counting sides of triangles gives
\[
3T=2E_{\mathrm{int}}+b=2E-b.
\]
Hence
\[
E=\frac{3T+b}{2}.
\]
Substituting this into Euler's formula yields
\[
(i+b)-\frac{3T+b}{2}+(T+1)=2,
\]
so
\[
T=2i+b-2.
\]
Therefore
\[
A=\frac{T}{2}=\frac{2i+b-2}{2}=i+\frac{b}{2}-1.
\]
This is Pick's theorem.
\end{proof}

\end{document}
